\section{Introduction}
\label{sec:introduction}

Tax-benefit microsimulation models translate statutory rules into
computable form over individual or household microdata, so that a policy
question -- who is eligible, what a household would owe or receive, what a
reform would cost -- can be answered directly rather than estimated from
aggregates. Established models of this kind carry a canonical citation that
describes the model once and is then cited by every downstream application:
\euromod{} has Sutherland and Figari's 2013 description
\citep{sutherland2013euromod}; \taxsim{} has Feenberg and Coutts's
introduction to the model.
\todo{The Feenberg-Coutts TAXSIM citation is named in this repository's seed
README but its bibliographic record has not been verified against a
primary source -- see docs/pending-citations.md for what to check before
adding it to bibliography/references.bib and citing it here with
\textbackslash citep\{feenberg1993taxsim\}.} \policyengineus{} does not yet
have this citation, despite implementing a comparable scope of federal,
state, and local rules: 67 registered programs across taxes and benefit
categories as of the version surveyed for this outline.
\evidence{policyengine\_us/programs.yaml; see Section~\ref{sec:coverage}}
This paper is that description for \policyengineus{}.

\subsection{Versioning}
\label{sec:versioning}

The paper is explicitly versioned. This is v1: it documents
\policyengineus{} as deployed in 2026, at the package version pinned in the
data-and-code-availability statement (Section~\ref{sec:disclosures}). A
model's rules content changes continuously -- \todo{cite release cadence,
e.g. commits/releases per year, once confirmed; the CHANGELOG shows patch
releases roughly weekly as of mid-2026} -- but that is ordinary maintenance,
not a change to the model's architecture, and does not by itself require a
new version of this paper. A v2 is warranted when the underlying computation
changes in a way that would make v1's description inaccurate rather than
merely dated: for example, a migration of the rules-encoding layer (tracked
separately under the PolicyEngine/Axiom rules-encoding effort) or a
wholesale replacement of the microdata foundation's estimation approach.
Absent such a migration, later editions of this paper should extend v1's
validation with updated numbers rather than restate its architecture.

\subsection{Scope and contributions}
\label{sec:scope}

This paper: (1) describes \policyengineus{}'s rules coverage, entity
structure, and parameter provenance (Section~\ref{sec:model}); (2) describes
the microdata foundation it runs over, the \populace{} stack, and how that
data is calibrated to administrative totals (Section~\ref{sec:data}); and
(3) validates the model's outputs against \taxsim{} cross-runs,
administrative aggregates, and published \cbo{}/\jct{}/\tpc{} scores where a
comparable estimate exists (Section~\ref{sec:validation}). Every claim in
Sections~\ref{sec:model}--\ref{sec:validation} is either backed by an
existing artifact in the \policyengineus{}, \populace{}, or \policybench{}
repositories, cited to its file or URL, or marked
\todo{evidence needed} where no existing artifact was found during outline
drafting -- this outline is deliberately explicit about which claims are
already assembled and which require new analysis before v1 can be
submitted.
